% BEGIN PROFESSOR FEEDBACK PF21: Results - RQ3
\section{RQ3: Distinct Motif Fingerprints per Category}

RQ3 evaluates category separability using motif fingerprints alone. It does
not compare fingerprints with conventional behavioural features for category
classification. The analysis contains 10{,}871 addresses across the categories listed in
Table~\ref{tab:scam_types}. Their distribution is very uneven: phishing accounts
for most addresses, while the seven smallest categories have only 25 to 47
addresses each in the test set.

Table~\ref{tab:res_rq3_global} reports the overall classification result, and
Table~\ref{tab:res_rq3_labels} gives the result for each category separately.
Figure~\ref{fig:res_rq3} compares the average fingerprints, and
Figure~\ref{fig:res_rq3_confusion} reports the confusion matrix.

The heatmap labels use layer prefixes for ordinary transfers (N), internal
transfers (I), ERC-20 (E20), ERC-721 (E721), and ERC-1155 (E1155).
The labels pre, mid, and post identify the three star classes, and two denotes
a two-node role. Direction codes use i for incoming and o for outgoing.
Category labels in the figures omit the shared suffix Scam.

\begin{figure}[p]
  \centering
  \ResGraphic{1.0}{rq3_fingerprint_heatmap}
  \caption{RQ3 category fingerprints using the globally selected two-week
  window and the default timestamp ordering. The two panels show the 40 local
  role components that vary most between categories in the development data,
  divided into two groups of 20 for readability. These components are selected
  from the five layers of 32 roles each, not from the 40-role HuggingFace
  catalogue.}
  \label{fig:res_rq3}
\end{figure}

\begin{figure}[p]
  \centering
  \ResGraphic{1.0}{rq3_confusion_matrix}
  \caption{RQ3 confusion matrix for the held-out Real-CATS addresses at the
  shared two-week window under the default timestamp ordering. Cells report
  numbers of addresses.}
  \label{fig:res_rq3_confusion}
\end{figure}

\IfFileExists{generated/results_rq3_global.tex}{\begin{table}[ht]
\centering
\caption{RQ3: overall separability of the criminal categories from their motif fingerprints (canonical ordering).}
\label{tab:res_rq3_global}
\small
\begin{tabular}{@{}l r@{}}
\toprule
\textbf{Quantity} & \textbf{Value} \\
\midrule
Global window & 5 min \\
Macro average precision & 0.333 \\
Macro F1 & 0.167 \\
Permutation $p$-value & 1.000 \\
Decision & not rejected \\
\bottomrule
\end{tabular}
\end{table}
}{\ResPendingTable{RQ3 overall category separability.}{tab:res_rq3_global}}
\IfFileExists{generated/results_rq3_labels.tex}{\begin{table}[ht]
\centering
\caption{RQ3: one-vs-rest average precision for each behaviour category at the globally selected window (canonical ordering). The address count and class share are included because average precision depends on class prevalence.}
\label{tab:res_rq3_labels}
\small
\begin{tabular}{@{}l r r r@{}}
\toprule
\textbf{Category} & \textbf{$n$} & \textbf{Share} & \textbf{AP} \\
\midrule
Phishing Scam & 7,525 & 0.692 & 0.756 \\
Impersonation Scam & 2,099 & 0.193 & 0.349 \\
Pigbutchering Scam & 131 & 0.012 & 0.191 \\
Contract Exploit Scam & 234 & 0.022 & 0.101 \\
Hack Scam & 220 & 0.020 & 0.092 \\
Fake Returns Scam & 154 & 0.014 & 0.071 \\
Rug Pull Scam & 179 & 0.016 & 0.053 \\
Blackmail Scam & 202 & 0.019 & 0.051 \\
NFT Airdrop Scam & 127 & 0.012 & 0.025 \\
\bottomrule
\end{tabular}
\end{table}
}{\ResPendingTable{RQ3 per-category detection.}{tab:res_rq3_labels}}

Using the default timestamp ordering, the development data selects a shared
window of two weeks. On the test set, macro AP is 0.1953 and macro F1 is 0.1914.
For comparison, the average class share across the nine one-vs-rest tasks is
$1/9=0.1111$. The permutation test gives $p<0.001$. It rejects the
label-permutation null
under all seven timestamp orderings, although their selected windows range from
one week to 30 days.

The one-vs-rest AP values range from 0.0244 for NFT Airdrop Scam to 0.7710 for
Phishing Scam. Every value is above the share of its category. Pigbutchering Scam reaches
0.2594 AP with a class share of only 0.0121. Phishing has the highest AP, but it
also represents 69.22\% of the sample. Table~\ref{tab:res_rq3_labels} reports
the class shares alongside AP.

A model based on log-transformed local role counts reaches macro AP 0.2058 and macro F1 0.1868,
compared with 0.1953 and 0.1914 for the composition-based fingerprint.
% END PROFESSOR FEEDBACK PF21

